%!TEX root = ../main.tex
\chapter{Related Work}
\label{ch:relatedwork}

\section{Adversarial Attacks on AMC Systems}

The vulnerability of deep-learning AMC classifiers to adversarial perturbations
was first systematically examined by Lin et al.~\cite{lin2020tactics}, who
applied FGSM and PGD to DNN-based modulation recognition, demonstrating that
single-step and iterative gradient attacks can degrade accuracy from above
90\% to below 10\%. Flowers et al.~\cite{flowers2019sigguard} provided a
comprehensive evaluation of adversarial evasion attacks in the wireless
context, showing that the threat is practical under realistic signal
conditions. Yang et al.~\cite{yang2018characterizing} characterized the
geometry of adversarial examples for wireless DNN classifiers, finding that
the decision boundary is significantly closer to clean samples in the IQ
domain than in equivalent image-classification tasks.

Bahramali et al.~\cite{bahramali2021robust} introduced over-the-air
adversarial attacks that survive the physical channel, revealing a more
stringent threat model where perturbations must remain effective after
transmission, fading, and noise. Silvaco et al.~\cite{silvaco2023adversarial_rf}
provided a comprehensive survey of adversarial attacks against AMC,
categorizing them by perturbation norm ($L_\infty$, $L_2$, $L_1$), attack
objective (untargeted vs.\ targeted), and knowledge assumption (white-box
vs.\ black-box). Zhang et al.~\cite{zhang2024stealthy} introduced stealthy
adversarial attacks that minimize spectral anomaly while maintaining high
misclassification rates, posing challenges for energy-based detectors. Kim
et al.~\cite{kim2024frequency} explored frequency-domain adversarial attacks
that concentrate perturbation energy in specific spectral bands, making
conventional band-limited filtering less effective.

Among attack methods, optimization-based attacks (CW~\cite{carlini2017towards},
EAD~\cite{chen2020ead}) consistently outperform gradient-based approaches
(FGSM~\cite{goodfellow2015fgsm}, PGD~\cite{madry2018pgd}) on IQ signals,
achieving near-zero accuracy against the AWN classifier at high SNR with
minimal perturbation norm. This is because IQ signals have low signal
amplitude ($\sim$0.02) relative to their dynamic range, allowing
optimization-based methods to find highly effective perturbations near the
classifier's decision boundary. In contrast to prior work that focuses solely
on attacks, we use these attacks as benchmarks to evaluate defense recovery.

\section{Defense Methods for AMC}

Defense strategies for AMC can be broadly divided into model-level and
input-level approaches. \emph{Model-level} defenses modify the classifier
through adversarial training~\cite{madry2018towards}, incorporating
adversarial examples during training to improve robustness. Tram\`{e}r et
al.~\cite{tramer2017ensemble} extended this with ensemble adversarial training
using transfers from multiple surrogate models. These methods are effective
against specific attacks but require retraining for each new threat, limiting
deployability in operational systems where the classifier is a fixed
checkpoint. Papernot et al.~\cite{papernot2016distillation} proposed defensive
distillation to smooth classifier gradients, reducing attack effectiveness,
but this approach has been shown to be bypassable by adaptive attacks.

\emph{Input-level} (input-transformation) defenses preprocess the input
before classification and require no model modification. Guo et
al.~\cite{guo2018input} systematically evaluated input transformations for
image adversarial defense, including JPEG compression, total variation
minimization, and image quilting, finding that FFT-based filtering is among
the most effective. Xu et al.~\cite{xu2018feature} proposed feature squeezing,
which reduces the feature space via color-depth reduction and spatial
smoothing, then uses disagreement between original and squeezed predictions
as a detection signal. Dziugaite et al.~\cite{dziugaite2016compression}
studied JPEG compression as a defense, showing that lossy compression removes
adversarial high-frequency components while preserving semantic content.
These image-domain insights motivate frequency-domain filtering for IQ
signals.

Cohen et al.~\cite{cohen2019certified} introduced randomized smoothing for
certified robustness: by taking a majority vote over classifier predictions
on Gaussian-noised copies of the input and computing a Clopper--Pearson
lower bound on the top-class probability $p_A$, the smoothed classifier
obtains a certified $L_2$ robustness radius $\sigma \Phi^{-1}(p_A)$. In our
evaluation we include an \emph{empirical randomized-smoothing baseline}
that follows the same majority-vote procedure with $k = 20$
Gaussian-noised copies at $\sigma = 0.01$ but does not compute the
Clopper--Pearson bound on $p_A$; the reported row therefore reflects the
empirical accuracy of the smoothed classifier rather than a certified
radius. Achieving a formal $L_2$ certificate would require hundreds to
thousands of noisy forward passes per sample~\cite{cohen2019certified},
which we treat as outside our target latency budget.

Recent work has applied frequency-domain ideas specifically to RF defense.
Kim et al.~\cite{kim2024frequency} proposed band-selective suppression based
on detecting spectral anomalies introduced by adversarial perturbations. Our
adaptive-$K$ method differs by operating on the full spectrum without prior
knowledge of which band the attack targets, making it attack-agnostic. To
the best of our knowledge, no prior work has performed a systematic
side-by-side comparison of frequency-domain adaptive methods and classical
communications signal-processing filters as adversarial defenses for AMC.

\section{Classical Signal-Processing Filters}

Classical signal processing offers a rich toolkit of denoising filters
developed for communications and estimation problems. The Wiener
filter~\cite{wiener1949extrapolation} is the minimum mean-square-error
(MMSE) estimator under the assumption of stationary signal and additive
noise, operating in the frequency domain by applying a magnitude-squared
SNR weighting to each spectral component. Kalman
filtering~\cite{kalman1960new} extends Wiener estimation to non-stationary
signals via a state-space model, recursively tracking signal evolution with
prediction and update steps. Both filters are well-suited to thermal noise
but were not designed for adversarial perturbations, which are structured
rather than random.

Savitzky-Golay smoothing~\cite{savitzky1964smoothing} fits a polynomial of
specified degree to a sliding window of samples using local least-squares,
producing a smoothing filter that preserves the shape of peaks and pulse
signals. Unlike frequency-domain filters, Savitzky-Golay operates in the
time domain and is particularly effective at preserving high-frequency signal
features while suppressing narrow noise spikes. Gaussian and FIR low-pass
filters are standard tools in digital communications for pulse shaping and
anti-aliasing; their regularization effect on IQ signals has been studied in
the context of matched filtering~\cite{proakis2008digital} and channel
estimation~\cite{haykin2002adaptive}, but their use as adversarial defenses
has not been examined in the AMC context.

The application of classical filters as adversarial defenses rests on the
hypothesis that adversarial perturbations behave like structured noise that
can be attenuated by denoising. However, this hypothesis is non-trivial:
unlike AWGN, adversarial perturbations are deliberately crafted to align with
the classifier gradient, and may concentrate energy in exactly the spectral
regions where the modulation signal is strongest. Our experiments test this
hypothesis directly by evaluating per-SNR-calibrated classical filters
against optimization-based attacks (CW, EAD) on the RML2016.10a dataset.

In contrast to the above works, we propose an adaptive spectral defense that
combines FFT Top-$K$ filtering with per-sample spectral-flatness routing,
and provide a systematic comparison of frequency-domain adaptive methods
against five per-SNR-calibrated classical filters and a Gaussian-noise
smoothing baseline as adversarial defenses for AMC. Our approach achieves
the best accuracy--latency trade-off among all evaluated methods in our
setting while remaining model-agnostic and suitable for low-latency
receiver-side deployment.
